\documentclass[11pt]{article}

\usepackage[margin=1in]{geometry}
\usepackage{amsmath, amssymb}
\usepackage{graphicx}
\usepackage{booktabs}
\usepackage{xcolor}
\usepackage{hyperref}
\usepackage{caption}
\usepackage{algorithm}
\usepackage{algpseudocode}

\newcommand{\TODO}[1]{\textcolor{red}{\textbf{[TODO: #1]}}}

\title{Bayesian and Neural Approaches to Word-Level Language Modeling:\\
A Comparative Study on the Sherlock Holmes Corpus}
\author{Devora Simi}
\date{\today}

\begin{document}
\maketitle

\begin{abstract}
We study next-word prediction on a corpus of Arthur Conan Doyle's Sherlock
Holmes novels and short stories as a controlled setting for comparing
statistical language models that make different commitments to Bayesian
inference. We implement and evaluate four model families on identical
train/validation/test splits: (i) an $n$-gram Markov model with Dirichlet
(additive) smoothing, whose smoothing constant is shown to be exactly the
Bayes estimator under a Dirichlet prior; (ii) a discrete-emission Hidden
Markov Model (HMM) fit by maximum-likelihood/MAP Baum-Welch EM, i.e.\ a
point-estimate latent-variable model; (iii) a fully Bayesian variant of the
same HMM fit by mean-field Variational Bayes EM (VB-EM), which maintains
Dirichlet \emph{posterior distributions} over the initial-state,
transition, and emission parameters rather than point estimates; and (iv)
an LSTM recurrent neural network language model, used both as a
point-estimate neural baseline and, via Monte Carlo Dropout, as an
approximate Bayesian neural model with predictive uncertainty. All models
are evaluated on held-out perplexity and top-$k$ next-word accuracy, the
latter reported both over all tokens and restricted to real words (i.e.\
excluding punctuation tokens). This report documents the problem
formulation, data pipeline, and modeling/evaluation methodology.
\emph{Experimental results are not yet included.}
\end{abstract}

\tableofcontents

% ============================================================
\section{Problem Statement}
% ============================================================

The task is \textbf{word-level next-token prediction}: given a prefix of
tokens $w_1, \dots, w_{t-1}$ from a sentence, predict a distribution over
the next token $w_t$ from a fixed vocabulary. This is the same underlying
predictive-text task that powers autocomplete and keyboard-suggestion
systems, and it is a convenient testbed for a broader question: \emph{how
much does the choice of Bayesian inference strategy matter for the same
modeling task and the same data?}

Rather than comparing unrelated model families, we deliberately structure
the comparison as three points on a single spectrum of how much
uncertainty a model represents over its own parameters:

\begin{enumerate}
    \item \textbf{Point estimation with a Bayesian smoothing prior.} The
    $n$-gram model and the standard HMM both use a Dirichlet prior only to
    regularize a single point estimate of their parameters (posterior mean
    or MAP), never propagating parameter uncertainty forward.
    \item \textbf{Full posterior inference.} The Variational Bayes HMM
    places the same conjugate Dirichlet priors on the same parameters as
    the point-estimate HMM, but maintains and updates full posterior
    distributions over them via mean-field VB-EM, so that predictions
    integrate over parameter uncertainty rather than plugging in a single
    value.
    \item \textbf{Approximate Bayesian deep learning.} The LSTM has no
    tractable parameter posterior, but Monte Carlo Dropout (Gal and
    Ghahramani, 2016) is used to approximate Bayesian model averaging over
    an implicit posterior over its weights, at prediction time.
\end{enumerate}

Holding the task, data, and evaluation protocol fixed lets us ask, for
each pair of models, whether the extra machinery of full or approximate
Bayesian inference actually changes predictive performance (perplexity,
top-$k$ accuracy) and/or yields useful predictive uncertainty, relative to
the corresponding point-estimate model.

A secondary, applied framing of the same task is a \textbf{predictive-text
/ autocomplete demo}: given a typed prefix, suggest the top-$k$ most
likely next words under each trained model (and, for the RNN, an
uncertainty estimate per suggestion via MC Dropout).

% ============================================================
\section{Data}
% ============================================================

\subsection{Source corpus}

The corpus is the complete Sherlock Holmes canon by Arthur Conan Doyle: nine public-domain works obtained from Project Gutenberg~\cite{doyle}, consisting of four novels and five short-story anthologies.

\begin{itemize}
    \item Novels: {A Study in Scarlet}, {The Sign of Four},
    {The Hound of the Baskervilles}, {The Valley of Fear}.
    \item Anthologies: {The Adventures of Sherlock Holmes},
    {The Memoirs of Sherlock Holmes}, {The Return of Sherlock
    Holmes}, {His Last Bow}, {The Case-Book of Sherlock Holmes}.
\end{itemize}

\subsection{Preprocessing pipeline}

Preprocessing (\texttt{src/preprocess.py}, \texttt{src/data\_utils.py})
follows the WikiText benchmark conventions~\cite{merity2016}, which use
the Moses tokenizer~\cite{koehn2007}:

\begin{enumerate}
    \item \textbf{Boilerplate removal.} Project Gutenberg license headers
    and footers are stripped.

    \item \textbf{Unit segmentation.} Anthologies are split into stories
    by locating each title from the ``Contents'' listing in the body text;
    novels stay whole. This gives 60 units: 56 stories and 4 novels.

    \item \textbf{Paragraphs.} Each sequence is one paragraph (text
    between blank lines), ending with \texttt{<eos>}; there is no sentence
    splitting.

    \item \textbf{Heading removal.} Titles, chapter headings, contents
    listings and bylines are dropped (WikiText keeps section titles; here
    they are layout, not content).

    \item \textbf{Normalization.} Variants of the same symbol are unified
    (e.g., curly/straight quotes).

    \item \textbf{Tokenization.} Case is kept, every punctuation mark is a
    token, and numbers are kept. Clitics are split (\texttt{don 't},
    \texttt{Holmes 's}), inner hyphens and number separators become markers
    (\texttt{well @-@ known}, \texttt{1 @,@ 000}), and title abbreviations
    stay whole (\texttt{Mr.}, \texttt{Dr.}).

    \item \textbf{Leak-free splitting.} Whole units are assigned to splits,
    so no story appears in two splits. Units are shuffled and
    filled into test, validation, then train, until each reaches its token
    target (10\%, 10\%, rest).

    \item \textbf{Vocabulary.} Built from the training split only: every
    token seen at least twice (WikiText uses three; this corpus is smaller so we used two), plus \texttt{<unk>} and
    \texttt{<eos>}.
\end{enumerate}

\subsection{Split statistics}

The training split contains 17{,}859 distinct token types, of which
11{,}106 occur at least twice; together with \texttt{<unk>} and
\texttt{<eos>} this gives a vocabulary of 11{,}108 entries, covering
98.8\% of training tokens (Table~\ref{tab:splits}). The out-of-vocabulary
rate is higher on the validation and test splits because they consist of
stories unseen in training, with their own character names and rare words.

\begin{table}[h]
\centering
\begin{tabular}{lrrrr}
\toprule
Split & Narrative units & Paragraphs & Tokens & \texttt{<unk>} rate \\
\midrule
Training & 47 & 11{,}208 & 583{,}089 & 1.16\% \\
Validation & 5 & 2{,}134 & 111{,}826 & 3.13\% \\
Test & 8 & 2{,}077 & 118{,}608 & 3.53\% \\
\bottomrule
\end{tabular}
\caption{Split sizes (60 narrative units in total; vocabulary of 11{,}108
entries, minimum training count 2).}
\label{tab:splits}
\end{table}

Test and validation units are listed below; all others are used for training:
\begin{itemize}
    \item \textbf{Test (8):} {A Study in Scarlet} (novel); from
    {Adventures}: The Copper Beeches; from {Memoirs}: The Reigate
    Squires; from {Return}: The Abbey Grange, The Missing
    Three-Quarter; from {Case-Book}: The Lion's Mane, The Retired
    Colourman, The Sussex Vampire.
    \item \textbf{Validation (5):} {The Hound of the Baskervilles}
    (novel); from {Adventures}: The Noble Bachelor; from
    {Memoirs}: The Cardboard Box, The ``Gloria Scott''; from
    {Case-Book}: The Creeping Man.
\end{itemize}

Paragraph lengths vary widely: the median paragraph has 29 tokens, the longest 710.

% ============================================================
\section{Models}
% ============================================================

All four model families are trained and evaluated on the same encoded
sentences (or, for the RNN, the same token stream) and share the same
notion of ``vocabulary,'' so their outputs are directly comparable.

\subsection{$n$-gram model with Dirichlet smoothing}

An order-$n$ Markov model estimates
\[
P(w_t \mid w_{t-n+1}, \dots, w_{t-1})
\]
from co-occurrence counts of each context and following word. Rather than
using raw maximum-likelihood counts (which assign zero probability to
unseen context/word pairs), the model uses additive (Dirichlet) smoothing
with concentration $\alpha$:
\[
P(w \mid \text{context}) = \frac{\text{count}(\text{context}, w) +
\alpha}{\text{count}(\text{context}) + \alpha V}, \qquad V = |\text{vocab}|.
\]
This is not merely a heuristic: it is exactly the posterior \emph{mean} of
a $\text{Dirichlet}(\alpha)$ prior placed over the categorical
distribution of ``next word given this context,'' and by
Dirichlet-multinomial conjugacy the posterior mean is also exactly the
Bayesian posterior-predictive probability for a single next draw. The
model is thus a genuine (if simple) Bayes estimator, using the same
conjugate-prior idea that underlies the HMM's smoothing below, applied
directly to $n$-gram counts. Implemented in \texttt{src/ngram.py}.

\textbf{Hyperparameters explored:} order $n$ (context length) and the
smoothing concentration $\alpha$.

\subsection{Hidden Markov Model (point estimate)}

A discrete-emission (categorical) HMM with $N$ latent states models a
sentence as a Markov chain of unobserved states
$z_1, \dots, z_T$ that emit the observed word sequence
$w_1, \dots, w_T$, parameterized by:
\begin{itemize}
    \item initial-state distribution $\pi \in \Delta^{N-1}$,
    \item transition matrix $A \in \mathbb{R}^{N \times N}$ (row-stochastic),
    \item emission matrix $B \in \mathbb{R}^{N \times V}$ (row-stochastic).
\end{itemize}
Parameters are fit by Baum-Welch (EM): the E-step computes state
posteriors $\gamma_t(i) = P(z_t = i \mid w_{1:T})$ and pairwise transition
posteriors via the scaled forward-backward algorithm; the M-step
re-estimates $\pi, A, B$ from these expected sufficient statistics, with a
small additive (Dirichlet-style) smoothing constant added to every count
to avoid degenerate zero probabilities. This yields \emph{point} (MAP)
estimates of $\pi, A, B$ -- a single number per parameter, with no
representation of estimation uncertainty. Training terminates on an
iteration cap or when the total data log-likelihood changes by less than
a tolerance. Implemented from scratch in \texttt{src/hmm.py} (exposing the
same attribute names as \texttt{hmmlearn} for interchangeability with the
Bayesian variant below).

\textbf{Hyperparameters explored:} number of latent states $N$.

\subsection{Variational Bayes HMM (full posterior)}

The Variational Bayes HMM (\texttt{src/vb\_hmm.py}) uses the identical
generative model as the point-estimate HMM above, but places explicit
Dirichlet priors on $\pi$, each row of $A$, and each row of $B$, and fits
the model by mean-field Variational Bayes EM (VB-EM) rather than plain
EM. Instead of point estimates, the model maintains Dirichlet
\emph{posterior} concentration parameters for each of $\pi, A, B$. The
E-step's forward-backward recursion plugs in the expected
log-probabilities under the current posterior,
\[
\widetilde{\theta} = \exp\big(\psi(\alpha) - \psi(\textstyle\sum_j \alpha_j)\big),
\]
where $\psi$ is the digamma function, rather than a point-estimate
probability -- the standard mean-field update for conjugate-exponential
models (Beal, 2003; Bishop, 2006, Ch.\ 10). By Jensen's inequality this
quantity is always $\le$ the naive ratio $\alpha_i / \sum_j \alpha_j$, a
``Bayesian shrinkage'' effect that makes VB-EM inherently more
conservative than MAP-EM about rare states and rare emissions. The M-step
is a conjugate update: each posterior is simply the prior plus expected
sufficient statistics, with no separate normalization needed. The
digamma function is implemented from scratch (shift-up recurrence plus
asymptotic expansion) to keep the model dependency-free.

The fitted model exposes \texttt{startprob\_}, \texttt{transmat\_}, and
\texttt{emissionprob\_} as the \emph{posterior mean} of each Dirichlet
(again, by conjugacy, exactly the Bayesian posterior-predictive
distribution for one more draw), making it a drop-in replacement for the
point-estimate HMM in every downstream evaluation and demo function. This
lets the study directly compare MAP-EM vs.\ VB-EM \emph{at matching
numbers of latent states}, isolating the effect of full posterior
inference from any confound of model capacity.

\textbf{Hyperparameters explored:} number of latent states $N$ (matched to
the point-estimate HMM sweep); Dirichlet prior concentrations
$\alpha_0, \beta_0, \pi_0$ for the transition, emission, and initial-state
priors.

\subsection{LSTM recurrent neural network}

A standard single-/multi-layer LSTM language model
(\texttt{src/rnn.py}): a learned word embedding feeds an LSTM, whose
hidden state at each position is projected by a linear decoder to logits
over the vocabulary, trained with cross-entropy loss via Adam, gradient
clipping, dropout regularization, and truncated backpropagation through
time (BPTT) over sentences concatenated into a single token stream
(delimited by \texttt{<eos>}). This is a point-estimate neural baseline:
a single set of network weights, no distribution over weights.

\textbf{Hyperparameters explored (grid sweep,
\texttt{src/experiment.py}):} hidden size, number of layers, dropout
rate; training uses early stopping on validation perplexity (patience-based),
restoring the best epoch's weights before final evaluation.

\textbf{Hyperparameter search (\texttt{src/tune\_rnn.py}):} embedding
size, dropout, and learning rate are additionally tuned with Optuna's
Tree-structured Parzen Estimator (TPE) sampler, at a \emph{fixed}
representative hidden size/layer count (searching architecture size
jointly would make the winning embedding size/dropout/learning rate only
valid for the specific architecture they happened to be paired with in a
given trial, which would not transfer safely to the separate
hidden-size/layer sweep above). The best configuration found is then
retrained for a longer budget with early stopping and full history
logging.

\subsubsection{Monte Carlo Dropout: approximate Bayesian inference}

Because the LSTM's exact weight posterior is intractable, we use Monte
Carlo Dropout (Gal and Ghahramani, 2016) as an approximate Bayesian
neural model. At prediction/evaluation time, dropout is kept
\emph{active} (rather than disabled, as is standard at test time), and
$M$ independent stochastic forward passes are run on the same input, each
with an independently sampled dropout mask. Averaging the resulting
softmax distributions approximates Bayesian model averaging over an
implicit posterior over the network's weights; the standard deviation
across samples for each candidate word gives an approximate predictive
uncertainty unavailable from a single deterministic forward pass. This is
used two ways:
\begin{itemize}
    \item \textbf{Demo-time uncertainty} (\texttt{predict\_demo.py}): top-$k$
    next-word suggestions annotated with a mean probability and standard
    deviation across MC samples.
    \item \textbf{Re-scored evaluation} (\texttt{evaluate\_rnn\_mc\_dropout},
    used in the sweep): perplexity and top-$k$ accuracy recomputed with MC
    Dropout averaging, to test whether Bayesian model averaging changes
    point predictive performance (not just uncertainty), and how many
    stochastic passes are needed before the averaged prediction stops
    changing.
\end{itemize}
Perplexity under MC Dropout is computed as $-\log$ of the \emph{mean}
predictive probability at the true token (the mathematically correct way
to score a Bayesian model average), not as the mean of each sample's
log-probability.

% ============================================================
\section{Evaluation Protocol}
% ============================================================

All four model families are scored with the same two metrics
(\texttt{src/evaluation.py}), computed identically (up to each model's
own notion of ``predict from the prefix seen so far'') across models:

\begin{itemize}
    \item \textbf{Perplexity:} $\exp\!\big(-\tfrac{1}{T}\sum_t \log
    P(w_t \mid w_{<t})\big)$ over all tokens in the evaluation split,
    including each sequence's first token (predicted from an empty/short
    context -- the $n$-gram model's shortest-available context, or the
    HMM's initial-state distribution).
    \item \textbf{Top-$k$ accuracy}, $k \in \{1, 5, 10\}$: whether the
    true next token is among the model's $k$ highest-probability
    predictions, computed incrementally (predict, then reveal the true
    token and update state) and \textbf{not} scored on each sequence's
    first token (no prior context to predict from).
\end{itemize}

Because punctuation tokens are highly predictable from local context (and
would inflate accuracy for any model that learns punctuation patterns
well), top-$k$ accuracy is reported \textbf{twice}: once over all tokens,
and once restricted to positions whose true next token is a real word
(i.e.\ excluding the punctuation subset of the vocabulary). The
words-only variant is the more meaningful measure of actual next-word
prediction quality.

All models are evaluated on the same validation and held-out test
splits described in Section~2.

% ============================================================
\section{Experimental Design}
% ============================================================

\texttt{src/experiment.py} runs a grid sweep over all four model
families on identical data splits, saving a checkpoint (model
weights/parameters, training config, and validation/test metrics) per
configuration for reproducibility:

\begin{itemize}
    \item \textbf{$n$-gram:} grid over order $\times$ smoothing $\alpha$.
    \item \textbf{HMM and VB-HMM:} grid over number of latent states,
    with the \emph{same} state-count grid used for both, enabling a
    direct MAP-EM vs.\ VB-EM comparison at matched capacity; VB-HMM prior
    concentrations held fixed.
    \item \textbf{LSTM:} grid over hidden size $\times$ number of layers
    $\times$ dropout, each trained with early stopping on validation
    perplexity; optionally, each trained configuration is additionally
    re-evaluated with MC Dropout at several sample counts to study
    convergence of the Bayesian model average.
\end{itemize}

Separately, \texttt{src/tune\_rnn.py} runs an Optuna TPE search over
embedding size, dropout, and learning rate at a fixed architecture size,
then retrains the best configuration with a longer budget and early
stopping.

Every configuration of every model is evaluated with the identical
protocol of Section~4 on the identical held-out test split, so that
perplexity and top-$k$ accuracy are directly comparable model-to-model
and not just within a model family.

% ============================================================
\section{Results}
% ============================================================

\TODO{This section will report, once experiments are finalized: (1) a
comparison table of best validation/test perplexity and top-1/top-5
accuracy (all-tokens and words-only) per model family, matching the
fields already logged to \texttt{results/comparison.csv}; (2)
per-family sweep plots -- perplexity vs.\ number of states (HMM,
VB-HMM), vs.\ order (n-gram, by $\alpha$), and vs.\ hidden size (RNN, by
layers/dropout); (3) the MC-Dropout sample-count convergence plot; (4)
the Optuna search plot and best hyperparameters found
(\texttt{results/optuna\_trials.csv}, \texttt{results/optuna\_search.png}); and
(5) a head-to-head summary figure of best-configuration perplexity and
accuracy across all four model families.}

\begin{table}[h]
\centering
\begin{tabular}{lrrrrr}
\toprule
Model & Config & Val.\ PPL & Test PPL & Top-1 acc.\ (words) & Top-5 acc.\ (words) \\
\midrule
$n$-gram & \TODO{best (order, $\alpha$)} & \TODO{} & \TODO{} & \TODO{} & \TODO{} \\
HMM & \TODO{best $N$} & \TODO{} & \TODO{} & \TODO{} & \TODO{} \\
VB-HMM & \TODO{best $N$} & \TODO{} & \TODO{} & \TODO{} & \TODO{} \\
LSTM & \TODO{best config} & \TODO{} & \TODO{} & \TODO{} & \TODO{} \\
LSTM + MC Dropout & \TODO{best config} & \TODO{} & \TODO{} & \TODO{} & \TODO{} \\
\bottomrule
\end{tabular}
\caption{\TODO{Best-configuration comparison across model families
(to be filled in from \texttt{results/comparison.csv}).}}
\end{table}

% \begin{figure}[h]
%   \centering
%   \includegraphics[width=0.8\textwidth]{figures/summary_comparison.png}
%   \caption{\TODO{Best test perplexity and words-only top-5 accuracy per model family.}}
% \end{figure}

% ============================================================
\section{Discussion and Conclusion}
% ============================================================

\TODO{To be written once results are available. This section should
address, at minimum: whether full-posterior VB-EM outperforms
point-estimate MAP-EM for the HMM at matched state counts (and whether
any gap narrows or widens with fewer latent states, where the
Bayesian-shrinkage effect of VB-EM should matter most); whether MC
Dropout changes the LSTM's point predictive performance or only adds
uncertainty estimates; and how the two Bayesian latent-variable models
(n-gram, HMM family) compare to the LSTM in the perplexity/accuracy
vs.\ model-complexity trade-off.}

% ============================================================
\begin{thebibliography}{9}

\bibitem{rabiner1989}
L.~R. Rabiner.
\newblock A tutorial on hidden Markov models and selected applications in speech recognition.
\newblock \emph{Proceedings of the IEEE}, 77(2):257--286, 1989.

\bibitem{beal2003}
M.~J. Beal.
\newblock \emph{Variational Algorithms for Approximate Bayesian Inference}.
\newblock PhD thesis, University College London, 2003.

\bibitem{bishop2006}
C.~M. Bishop.
\newblock \emph{Pattern Recognition and Machine Learning}.
\newblock Springer, 2006. Chapter 10: Approximate Inference.

\bibitem{gal2016}
Y.~Gal and Z.~Ghahramani.
\newblock Dropout as a Bayesian approximation: Representing model uncertainty in deep learning.
\newblock In \emph{Proceedings of the 33rd International Conference on Machine Learning (ICML)}, 2016.

\bibitem{akiba2019}
T.~Akiba, S.~Sano, T.~Yanase, T.~Ohta, and M.~Koyama.
\newblock Optuna: A next-generation hyperparameter optimization framework.
\newblock In \emph{Proceedings of the 25th ACM SIGKDD International Conference on Knowledge Discovery \& Data Mining}, 2019.

\bibitem{doyle}
A.~C. Doyle.
\newblock The Sherlock Holmes canon: \emph{A Study in Scarlet}, \emph{The Sign of Four},
\emph{The Adventures of Sherlock Holmes}, \emph{The Memoirs of Sherlock Holmes},
\emph{The Hound of the Baskervilles}, \emph{The Return of Sherlock Holmes},
\emph{The Valley of Fear}, \emph{His Last Bow}, \emph{The Case-Book of Sherlock Holmes}.
\newblock Project Gutenberg, \url{https://www.gutenberg.org/}.

\bibitem{merity2016}
S.~Merity, C.~Xiong, J.~Bradbury, and R.~Socher.
\newblock Pointer sentinel mixture models.
\newblock \emph{arXiv preprint arXiv:1609.07843}, 2016.
(Introduces the WikiText-2 and WikiText-103 datasets.)

\bibitem{koehn2007}
P.~Koehn, H.~Hoang, A.~Birch, C.~Callison-Burch, M.~Federico, N.~Bertoldi, B.~Cowan,
W.~Shen, C.~Moran, R.~Zens, C.~Dyer, O.~Bojar, A.~Constantin, and E.~Herbst.
\newblock Moses: Open source toolkit for statistical machine translation.
\newblock In \emph{Proceedings of the 45th Annual Meeting of the ACL, Demo and Poster Sessions}, 2007.

\end{thebibliography}

\end{document}
